\indent ಇದೇ ರೀತಿಯಾಗಿ ಅರಿಶಿನ ಕುಂಕುಮ, ಫಲ ಪುಷ್ಪ, ಗಿಡ ಮರ ಬಳ್ಳಿ ಕೈಕಾಲು ಮುಖ, ಸೂರ್ಯ ಚಂದ್ರ ನಕ್ಷತ್ರ, ಗುಡುಗು ಸಿಡಿಲು ಮಿಂಚು ಶಬ್ದ ಸ್ಪರ್ಶ ರೂಪ ರಸ ಗಂಧ-ದ್ವಂದ್ವ ಸಮಾಸ ಪದಗಳು.\par

\kbvsechead{(೬)\hspace{0.45em}ಬಹುವ್ರೀಹಿ ಸಮಾಸ :}

\indent “ಹಣೆಗಣ್ಣ” ಎಂಬ ಪದವನ್ನು ಹಣೆಯಲ್ಲಿ + ಕಣ್ಣುಳ್ಳವನು ಎಂದು\par

\indent “ಫಾಲನೇತ್ರ” ಎಂಬ ಪದವನ್ನು ಫಾಲದಲ್ಲಿ + ನೇತ್ರವುಳ್ಳವನು ಎಂದು ಬಿಡಿಸಿ ಬರೆಯಬಹುದು.\par

\indent ಹಣೆಯಲ್ಲಿ + ಕಣ್ಣುಳ್ಳವನು = ಹಣೆಗಣ್ಣ\par

\indent ಈ ಉದಾಹರಣೆಯಲ್ಲಿ ಪೂರ್ವಪದದ ಅರ್ಥವಾಗಲೀ ಉತ್ತರಪದದ ಅರ್ಥವಾಗಲೀ ಪ್ರಧಾನವಲ್ಲ; ಹಣೆಯಲ್ಲಿ ಕಣ್ಣುಳ್ಳವನು ಯಾವನೋ ಅವನು ಪ್ರಧಾನ. ಇದರಿಂದ ಸಮಾಸದ ಯಾವುದೊಂದು ಪದವೂ ಪ್ರಧಾನವಾಗಿ ಅನ್ಯಪದದ ಅರ್ಥ ಪ್ರಧಾನವಾಗಿರುವುದೆಂದು ತಿಳಿದು ಬರುವುದು. ಲಕ್ಷಣವುಳ್ಳ ಸಮಾಸಕ್ಕೆ ಬಹುವ್ರೀಹಿ ಸಮಾಸವೆಂದು ಹೆಸರು.\par

\begin{kbvrule}
\noindent{\kbvsecfont ವ್ಯಾಖ್ಯೆ :}\hspace{0.5em}ಎರಡು ನಾಮಪದಗಳು ಸೇರಿ ಸಮಾಸವಾದಾಗ ಆ ಪದದ ಅರ್ಥವು ಪ್ರಧಾನವಾಗಿ ಉಳ್ಳ ಸಮಾಸಕ್ಕೆ ಬಹುವ್ರೀಹಿ ಸಮಾಸವೆಂದು.
\end{kbvrule}

\vspace{3pt plus 2pt}\noindent\hspace*{1.2em}\kbvlabel{ಉದಾ}\par\vspace{1pt}

\begin{kbvderiv}
\vspace{4pt plus 3pt minus 1pt}
\begin{center}\begin{adjustbox}{max width=\textwidth}
{\renewcommand{\arraystretch}{1.22}\setlength{\tabcolsep}{4pt}
\begin{tabular}{l c l}
ಒಂದು + ಕಣ್ಣುಳ್ಳವ & {\color{kbvgrey}$=$} & \kbvresult{ಒಕ್ಕಣ್ಣ} \\
ನಾಲ್ಕು + ಮೊಗವುಳ್ಳವ & {\color{kbvgrey}$=$} & \kbvresult{ನಾಲ್ಮೊಗ} \\
ಕೆಂಪು + ಕಣ್ಣುಳ್ಳವ & {\color{kbvgrey}$=$} & \kbvresult{ಕೆಂಗಣ್ಣ} \\
ಚಂದ್ರನಂತೆ + ಮುಖವುಳ್ಳವಳು & {\color{kbvgrey}$=$} & \kbvresult{ಚಂದ್ರಮುಖಿ} \\
ವಜ್ರದಂತೆ + ಶರೀರವುಳ್ಳವ & {\color{kbvgrey}$=$} & \kbvresult{ವಜ್ರಶರೀರಿ} \\
\end{tabular}}
\end{adjustbox}\end{center}
\vspace{4pt plus 3pt minus 1pt}
\end{kbvderiv}

\kbvsechead{(೭)\hspace{0.45em}ಕ್ರಿಯಾ ಸಮಾಸ:}

\begin{kbvrule}
\noindent{\kbvsecfont ವ್ಯಾಖ್ಯೆ :}\hspace{0.5em}ಪೂರ್ವಪದವು ಸಾಮಾನ್ಯವಾಗಿ ದ್ವಿತೀಯಾ ವಿಭಕ್ತಿ ನಾಮಪದವಾಗಿದ್ದು ಕ್ರಿಯಾಸೂಚಕ ಉತ್ತರಪದದೊಂದಿಗೆ ಸೇರಿ ಆಗುವ ಸಮಾಸಕ್ಕೆ ಕ್ರಿಯಾ ಸಮಾಸವೆಂದು ಹೆಸರು.
\end{kbvrule}
